\section{Conclusion}
\label{sec:conclusion}

This study developed and evaluated a modular finite-volume solver for the one-dimensional linear advection equation with a prescribed time-dependent inlet condition. Second-order upwind and second-order central spatial discretizations were considered in combination with the explicit midpoint Runge--Kutta method (RK2) and the classical fourth-order Runge--Kutta method (RK4). The analytical solution was used as a reference to assess numerical accuracy, while grid refinement, time-step refinement, and stability analysis were used to investigate the behavior of the numerical schemes.

For the second-order upwind scheme combined with RK2, grid refinement at a constant CFL number of $0.25$ produced an observed convergence order of approximately $1.94$, consistent with the expected second-order accuracy. Time-step refinement on a fixed grid further showed that the RK2 error approached the spatial discretization error as the time step decreased. With RK4, the temporal error was already sufficiently small for the spatial error to dominate. Consequently, grid refinement at a constant CFL number yielded an observed order of approximately $1.94$ for both time integrators. These results demonstrate that increasing the temporal order does not necessarily increase the overall accuracy when the spatial discretization is the dominant source of error.

The Fourier stability analysis of the second-order upwind discretization gave a theoretical CFL limit of $C\leq 0.5$ for RK2. Tests using the full solver were consistent with this prediction: the solution remained bounded at CFL numbers of $0.4$ and $0.5$, whereas instability was observed at approximately $0.571$. The result at the limiting value should nevertheless be interpreted as marginal stability rather than as a robust operating condition. For RK4, the corresponding theoretical limit was approximately $C_{\max}=0.6963$, allowing a larger time step than RK2 for the same spatial grid. A numerical search yielded a slightly lower estimate because of its finite resolution.

The stability analysis also showed that the second-order central discretization combined with RK2 is unstable for nonzero Fourier modes at any positive CFL number. Therefore, a small time step or apparently acceptable behavior over a short simulation does not establish the long-time stability of this combination.

Overall, the second-order upwind scheme provides a suitable baseline for the present advection problem. RK2 offers a straightforward and effective time integrator when used within its stability limit, whereas RK4 provides a wider stability range and smaller temporal error. In either case, the accuracy of the complete solver remains constrained by the second-order spatial discretization. Finally, the Fourier analysis characterizes the stability of the interior discretization under periodic-mode assumptions; confirming stability for the non-periodic inlet-boundary problem requires consideration of the actual boundary treatment as well.
